\section{Experiments and Results}
\label{sec:experiments_results}

The experimental evaluation was centred on the main empirical constraint of the cohort: multimodal evidence was available at scale, whereas report-level semantic supervision was incomplete and unevenly distributed across centres. This setting provides a direct test of LCAD--RASA, which uses structured report semantics as an alignment signal without assuming that physician-authored reports are available for every patient. The results are presented as a connected evidence chain. We first characterise the supervision bottleneck, then evaluate whether pseudo reports provide modality-grounded semantic targets, examine whether these targets support section-wise multimodal alignment, and assess their value under report scarcity. We then place the held-out diagnostic results against both internal variants and same-split external baselines, before using ablation, perturbation, cross-centre, safety, and scalability analyses to define the scope and limitations of the framework.

\subsection{Five-centre cohort characteristics and report-supervision imbalance}
\label{subsec:cohort_results}

The final analytic cohort comprised 1,897 multimodal cervical screening cases from five centres and included 137,591 evaluable OCT and colposcopy images (Table~\ref{tab:cohort_summary}, ``Cohort scale and case-level report supervision''; Figure~\ref{fig:centre_supervision}, ``Centre-level cohort scale and report-supervision imbalance''). This scale was sufficient for multimodal representation learning, but it was not matched by complete report-level supervision. Archived clinician-authored reports were available for 744 cases, whereas 1,153 cases had complete multimodal evidence but no archived report and were therefore eligible for pseudo-report completion.

The imbalance was also highly centre-dependent. Enshi retained real reports for all 406 cases, Jingzhou retained reports for 334 of 406 cases, Xiangyang retained only 4 reports among 500 cases, and Wuhan and Shiyan had no centre-wide real-report archives. The resulting distribution indicates that the principal limitation was not image acquisition, but the uneven availability of semantic supervision. This observation motivates a report-aware strategy that can use real reports when they exist while extending structured semantic alignment to cases in which multimodal evidence is present but report records are missing.

\subsection{Modality-grounded pseudo-report generation across sources}
\label{subsec:pseudo_report_results}

Given this supervision gap, the first methodological requirement was to determine whether pseudo reports could provide meaningful semantic augmentation rather than merely converting labels into text. The label-template baseline served as a negative control. Although it satisfied the expected report schema, it had no modality support, with a mean modality-support rate of 0.000, and showed marked textual repetition, with a maximum duplicate fraction of 0.867 (Table~\ref{tab:pseudo_source_comparison}, ``Structured pseudo-report source comparison''; Figure~\ref{fig:pseudo_source_comparison}, ``Structured pseudo-report source comparison across template, rule-based, and local LLM sources''). Thus, schema validity alone did not establish useful supervision: a report could be structurally complete while still failing to encode OCT, colposcopy, or clinical evidence.

Rule-based and local embedding-assisted LLM pseudo reports showed a different pattern. Both preserved OCT, colposcopy, and clinical section support, with a mean modality-support rate of 1.000. The rule-based baseline reached a label-consistency mean of 0.628, and the local LLM baseline reached the same mean while retaining modality-grounded sections. These results indicate that structured pseudo reports can supply section-level semantic anchors linked to the observed modalities, rather than simply restating diagnostic labels in a report-like format.

The external-provider comparison examined whether this pseudo-reporting mechanism was tied to a single LLM source. Template, rule-based, local embedding LLM, Qwen-Plus, GLM-4.7-Flash, Gemini-3.1-Pro, and GPT-5.5 were evaluated under a common structured-report schema (Table~\ref{tab:llm_provider_comparison}, ``LLM provider comparison for structured pseudo-report generation''; Figure~\ref{fig:api_provider_summary_heatmap}, ``LLM provider comparison for structured pseudo-report generation''; Figure~\ref{fig:api_quality_heatmap}, ``Stage 1 provider quality metrics and risk-oriented quality indicators''; Figure~\ref{fig:api_reliability}, ``Provider latency, modality support, and generation reliability''). All evaluated sources produced schema-valid structured outputs in their available cohorts. In the extended 100-case cohort, GPT-5.5 generated 100 valid structured reports, with a schema-valid rate of 1.000, section completeness of 0.810, modality support of 0.810, contradiction rate of 0.050, hallucination rate of 0.000, mean latency of 11.39 seconds per case, and an estimated cost of US\$27.73 per 1,000 cases. Qwen-Plus, GLM-4.7-Flash, and Gemini-3.1-Pro were retained as pilot cohorts, and their results are therefore interpreted as feasibility evidence rather than as a definitive ranking of providers.

All external API evaluations used only de-identified structured evidence. Raw images, image paths, patient names, internal patient identifiers, and hospital identifiers were not transmitted. This design keeps pseudo-report generation within a controlled learning pipeline and preserves the distinction between API-generated structured text and physician-authored clinical documentation.

\subsection{Section-wise semantic alignment of multimodal representations}
\label{subsec:alignment_results}

The modality-grounded structure of the pseudo reports enabled a direct test of whether report sections could serve as alignment targets for multimodal representations. OCT embeddings were matched to \texttt{oct\_findings}, colposcopy embeddings to \texttt{colposcopy\_findings}, clinical embeddings to \texttt{clinical\_context}, and fused representations to \texttt{impression}. Full LCAD--RASA achieved a macro mean reciprocal rank (MRR) of 0.051 across the four sections. Removing section alignment reduced macro MRR to 0.021, while simple concatenation fusion remained similarly low at 0.022 (Table~\ref{tab:theme1_alignment_ablation}, ``Direct RASA alignment ablation using modality-section retrieval''; Figure~\ref{fig:alignment_retrieval_mrr}, ``Modality-section retrieval MRR for Theme 1 cross-modal semantic alignment'').

The real-report-only variant achieved the highest retrieval-only macro MRR, at 0.057, which is consistent with clinician-authored reports acting as strong semantic targets when available. However, this retrieval advantage does not resolve the coverage problem created by sparse and centre-imbalanced report archives. LCAD--RASA is therefore not intended to replace real reports; rather, it extends report-guided alignment to cases with complete multimodal evidence but missing report supervision, while maintaining a clear distinction between real-report and pseudo-report sources.

The staged API alignment analysis further showed that section-level alignment remained measurable across generation sources. Gemini-3.1-Pro achieved the highest available pilot macro MRR of 0.366, whereas GPT-5.5 achieved a macro MRR of 0.063 on the extended cohort (Figure~\ref{fig:api_alignment_mrr}, ``Stage 2 provider alignment analysis: macro MRR and section-specific MRR''; Figure~\ref{fig:api_section_mrr}, ``Provider-level section MRR across report sections''). Because cohort size and evaluation stage differed across providers, these results should not be interpreted as a definitive provider comparison. Their main implication is that structured report sections provide tractable alignment targets whose behaviour can be quantified across pseudo-report sources.

\subsection{Pseudo-report augmentation under report scarcity}
\label{subsec:scarcity_results}

The practical value of pseudo-report augmentation became most apparent when real-report supervision was deliberately restricted. In the report-supervision scarcity curve, the largest improvement occurred in the lowest-supervision regime (Table~\ref{tab:scarcity_curve}, ``Report-supervision scarcity curve using a lightweight downstream surrogate''; Figure~\ref{fig:scarcity_curve}, ``Report-supervision scarcity curve comparing real-report-only and pseudo-report-augmented surrogates''). At 10\% real-report availability, the LCAD-augmented surrogate achieved AUROC 0.808 and F1 0.500, compared with AUROC 0.630 and F1 0.316 for the real-report-only surrogate. At full real-report availability, the augmented surrogate also remained higher, with AUROC 0.838 and F1 0.542 compared with AUROC 0.799 and F1 0.471 for the real-report-only surrogate.

This pattern aligns with the cohort-level bottleneck identified above. Pseudo-report augmentation was most useful in the setting where images and structured clinical variables were available but semantic report supervision was incomplete, rather than functioning merely as an increase in textual volume. The staged API downstream surrogate was intentionally conservative: GPT-5.5 achieved AUROC 0.630 and F1 0.229 at 10\% real-report availability on the extended cohort, whereas Gemini-3.1-Pro achieved AUROC 0.725 and F1 0.379 in its pilot cohort (Figure~\ref{fig:api_scarcity_auc}, ``Stage 3 downstream scarcity surrogate for selected API-generated pseudo reports''). These analyses support the role of structured pseudo-report augmentation under report scarcity while remaining complementary to full end-to-end retraining evidence.

\subsection{Held-out diagnostic performance of report-aware variants}
\label{subsec:diagnostic_results}

The semantic-alignment analyses were then evaluated in terms of diagnostic prediction under the locked retrospective protocol. On the held-out test set ($n=288$), Full LCAD--RASA achieved AUROC 0.832 (95\% CI 0.757--0.897) and F1 0.611 (95\% CI 0.458--0.737) at the validation-selected threshold of 0.37 (Table~\ref{tab:main_comparison}, ``Held-out test performance at validation-selected thresholds''; Figure~\ref{fig:main_auc}, ``Held-out AUROC with bootstrap confidence intervals''; Figure~\ref{fig:main_metrics}, ``Held-out multi-metric profile across model variants''; Figure~\ref{fig:main_auc_f1}, ``AUROC--F1 trade-off across main model variants''). Pseudo-augmented LCAD training achieved AUROC 0.826 (95\% CI 0.747--0.894) and F1 0.545 (95\% CI 0.395--0.674). Real-report-only training achieved AUROC 0.725, whereas simple concatenation fusion achieved AUROC 0.472.

The ordering of internal variants was consistent with the preceding evidence. Generic concatenation was inadequate for the multimodal setting; real-report-only training was limited by incomplete report coverage; pseudo-augmented learning improved over real-report-only learning; and the full LCAD--RASA configuration achieved the highest point estimates among the report-aware internal variants. Corrected paired bootstrap rechecking supported clear AUROC differences between Full LCAD--RASA and weaker internal comparators such as simple concatenation fusion and real-report-only training, but not between Full LCAD--RASA and pseudo-augmented LCAD. The held-out results therefore support favourable internal performance under the locked retrospective protocol, without establishing clinical superiority.

\subsection{Same-split comparison with external multimodal baselines}
\label{subsec:external_baseline_results}

To place the internal gains in a broader comparative context, we added a same-split external baseline block using the identical train/validation/test split, validation-selected threshold rule, held-out test set, and bootstrap confidence-interval protocol. The comparison included clinical-only logistic regression, clinical-only gradient boosting, OCT-only and colposcopy-only frozen-embedding MLP classifiers, a late-fusion MLP, a cross-attention multimodal transformer, and a CLIP-style contrastive multimodal baseline without report-section supervision (Table~\ref{tab:external_baselines}, ``Same-split external baseline comparison on the locked held-out test set''; Figure~\ref{fig:external_baselines_auc}, ``Same-split external baseline AUROC comparison with bootstrap confidence intervals''; Figure~\ref{fig:external_baselines_metrics}, ``Metric profile for Full LCAD--RASA and external baselines under the locked held-out protocol''). Clinical-only baselines used age, HPV, and TCT only; centre identifiers, report text, pathology-like attributes, and labels were excluded from their inputs.

Full LCAD--RASA was competitive with several conventional baselines, including clinical-only logistic regression (AUROC 0.830), colposcopy-only embedding MLP (AUROC 0.827), and late-fusion MLP (AUROC 0.824). It did not, however, dominate all stronger external comparators. Clinical-only HistGradientBoosting reached AUROC 0.862, the cross-attention multimodal transformer reached AUROC 0.847, and the CLIP-style contrastive multimodal baseline reached AUROC 0.888. In the corrected paired bootstrap recheck, the contrastive baseline exceeded Full LCAD--RASA in AUROC by 0.056 (95\% bootstrap interval 0.024--0.092 in favour of the comparator; two-sided $p=0.002$), whereas differences against clinical-only logistic regression, clinical-only HistGradientBoosting, late-fusion MLP, and cross-attention transformer were not statistically conclusive (Table~\ref{tab:paired_bootstrap_recheck}, ``Corrected paired bootstrap AUROC recheck''; Figure~\ref{fig:external_pairwise_recheck}, ``Corrected paired bootstrap AUROC differences between Full LCAD--RASA and external baselines'').

These results refine the interpretation of LCAD--RASA. The framework should not be presented as a universally superior discriminative classifier when AUROC is the only criterion. Its primary contribution lies in report-aware multimodal learning: it converts incomplete report supervision into structured, modality-grounded semantic anchors; provides section-level alignment targets; and supports auditable pseudo-report generation under sparse report supervision. The strong performance of the contrastive baseline suggests a complementary direction for future development, in which report-section supervision is integrated with contrastive multimodal pretraining.

\subsection{Contribution and limitations of report-aware alignment components}
\label{subsec:ablation_results}

Ablation analyses were used to examine whether the observed performance patterns were coherent with the report-aware alignment design. Removing section alignment reduced direct alignment quality, with macro MRR decreasing from 0.051 for Full LCAD--RASA to 0.021 (Table~\ref{tab:theme1_alignment_ablation}, ``Direct RASA alignment ablation using modality-section retrieval''). In the held-out comparison, the LCAD variant without section alignment retained a similar AUROC of 0.826, but its F1 score decreased to 0.510 compared with 0.611 for Full LCAD--RASA (Table~\ref{tab:main_comparison}, ``Held-out test performance at validation-selected thresholds''). This divergence suggests that section-level alignment may influence threshold-dependent error balance and operating-point behaviour even when rank-based discrimination is similar.

Additional modality, quality-control, alignment-weight, and component ablations are reported in Supplementary Tables~S1 and S3--S5 and Figure~\ref{fig:rasa_lambda} (``Alignment-weight sweep for the RASA objective''), Figure~\ref{fig:rasa_component_boxenplot} (``Modality ablation and RASA component ablation''), and Figure~\ref{fig:qc_ablation} (``LCAD quality-control ablation''). These analyses are best interpreted as mechanism diagnostics and internal consistency checks rather than as final-model superiority tests, because several variants were run under fixed quick-budget settings and some produced point estimates close to the full model. In particular, the QC ablation did not materially change downstream AUROC/F1 under the current retrospective configuration. QC is therefore treated as a procedural safeguard for pseudo-report generation, not as a major independent driver of diagnostic performance. Overall, the ablations support the internal coherence of the RASA design while avoiding a causal interpretation of any single component.

\subsection{Perturbation-based evidence for modality-grounded decoding}
\label{subsec:perturbation_results}

Perturbation analysis examined whether generated report sections were sensitive to the modality evidence from which they were expected to derive. Removal or disruption of a modality produced the largest degradation in the matched report section (Table~\ref{tab:perturbation_matrix}, ``Upgraded perturbation sensitivity matrix''; Figure~\ref{fig:perturbation_heatmap}, ``Perturbation condition by report-section similarity matrix''; Figure~\ref{fig:perturbation_lineplot}, ``Section-level degradation profiles under modality perturbations''; Figure~\ref{fig:risk_delta}, ``Risk-score shifts under report-generation perturbations''; Figure~\ref{fig:theme1_perturbation}, ``Theme 1 upgraded perturbation sensitivity matrix''). Masking OCT produced a 0.743 drop in the OCT findings section, masking colposcopy produced a 1.000 drop in the colposcopy findings section, and masking clinical instruction produced a 0.833 drop in the clinical-context section.

Label-only inference produced the largest overall report drop of 0.466 and the largest absolute risk shift of 0.371. The degradation profile indicates that report generation was not fully recoverable from class labels alone and that section content remained sensitive to modality-specific evidence. This supports the interpretation that LCAD--RASA performs modality-grounded decoding rather than fixed-template schema filling. The evidence remains perturbational rather than causal; generated reports should therefore not be interpreted as causal clinical explanations.

\subsection{Reference-stratified performance and centre-level generalisation}
\label{subsec:loco_results}

Reference-stratified testing assessed whether held-out performance was confined to cases with archived reports. Full LCAD--RASA achieved AUROC 0.824 among cases with reference reports ($n=108$) and AUROC 0.850 among cases without reference reports ($n=180$). Comparable performance across these strata suggests that the model's held-out behaviour was not restricted to the real-report subset, although the analysis remains retrospective and conditioned on the available cohort composition.

Cross-centre evaluation imposed a stricter test of generalisation. Under strict leave-one-centre-out retraining with the fixed quick-budget protocol, held-out AUROC was 0.702 for Enshi, 0.648 for Jingzhou, and 0.382 for Xiangyang (Supplementary Tables~S2 and S2b; Figure~\ref{fig:loco_heatmap}, ``Leave-one-centre-out evaluation: AUROC heatmap and forest-style centre summary''; Figure~\ref{fig:loco_forest}, ``Strict LOCO centre-wise behaviour under fixed quick-budget retraining''). Eval-only LOCO using a global checkpoint is reported separately and is not pooled with strict LOCO retraining. These results identify centre shift as a material limitation of the current model. The evidence supports report-aware fusion in retrospective held-out testing, but it does not justify an unrestricted external-validation claim across centres.

\subsection{Robustness, safety, and computational feasibility}
\label{subsec:robustness_safety_scalability}

The final analyses evaluated whether the pipeline remained reproducible and computationally tractable at cohort scale. The supplementary evaluation included masking validation, modality ablations, quality-control ablations, multi-seed stability, report-level safety metrics, and runtime summaries (Supplementary Tables~S7--S11; Supplementary Figure~\ref{fig:supplementary_robustness}, ``Supplementary robustness checks: masking validation and multi-seed stability''). The pipeline processed approximately 137k images, and the publishable model checkpoints were approximately 80 MB. These results support the retrospective feasibility of large-scale multimodal analytics with LCAD--RASA. At the same time, the external baseline and centre-shift findings indicate that prospective validation, stronger domain-shift mitigation, and integration with contrastive multimodal pretraining remain necessary next steps.
